\documentclass{article}

\usepackage[letterpaper,margin=1in]{geometry}
\usepackage{amsmath}
\usepackage{amssymb}
\usepackage{amsfonts}
\usepackage{mathtools}
\usepackage{graphicx}
\usepackage{booktabs}
\usepackage{hyperref}
\usepackage{float}
\usepackage{caption}
\usepackage{subcaption}
\usepackage{bm}

\hypersetup{
    colorlinks=false,
    pdfborder={0 0 0}
}

% Custom commands for consistency
\newcommand{\Rgeom}{R_{\mathrm{geom}}}
\newcommand{\Rrec}{R_{\mathrm{rec}}}
\newcommand{\Rtot}{R_{\mathrm{tot}}}
\newcommand{\Rdiff}{R_{\mathrm{diff}}}
\newcommand{\Rrxn}{R_{\mathrm{rxn}}}
\newcommand{\Jmax}{J_{\mathrm{max}}}
\newcommand{\kastar}{k_a^{*}}
\newcommand{\Nrec}{N_{\mathrm{rec}}}
\newcommand{\Nsurf}{N_{\mathrm{surf}}}
\newcommand{\Acyl}{A_{\mathrm{cylinder}}}

\title{Microscale Chemoreception in a Cylinder--Cylinder Model}

\begin{document}
\maketitle
\begin{abstract}
We investigate how cells in human tissues sense hormones, specifically insulin, through chemoreception.
This process spans multiple spatial scales, from whole-body insulin production and clearance to the microscopic dynamics of individual cells and their receptors.
In this work, we focus on the microscale and model a simplified system in which tissue cells surround a single blood capillary.
Our aim is to understand how geometry, receptor number, local insulin concentration, and binding--unbinding kinetics together determine the flux of insulin into cells, and whether this sensing operates near an optimal regime.
\end{abstract}

\section{Model geometry and approach outline}\label{sec:geometry}
To model diffusion of insulin from capillary to neighboring cells,
we introduce a geometry of two concentric cylinders,
with a constant ligand concentration imposed at the inner cylinder of radius $R_1$ and absorbing insulin receptors distributed on the outer cylinder of radius $R_2$ (figure~\ref{fig:capillary_cells}).
We specify the associated boundary conditions and replace the heterogeneous receptor distribution with an effective homogenized boundary condition,
enabling analytical treatment of the diffusion problem.
We then incorporate receptor binding and unbinding kinetics and, finally, validate the resulting analytical predictions by comparison with stochastic random-walk simulations.

\begin{figure}[H]
    \centering
    % Placeholder for Figure 1: 3D representation of the capillary and surrounding cells.
    \fbox{\parbox[c][6cm][c]{0.8\textwidth}{\centering \textit{Placeholder for Figure 1: 3D representation of the capillary and surrounding cells.}}}
    \caption{\textbf{3D representation of the capillary and surrounding cells.} The orange disks denote receptors uniformly distributed along the cell membrane.}
    \label{fig:capillary_cells}
\end{figure}

\begin{table}[H]
    \centering
    \caption{\textbf{System variables and parameters used in the model.}}
    \label{tab:parameters}
    \begin{tabular}{@{}lll@{}}
        \toprule
        \textbf{Variable} & \textbf{Value} & \textbf{Description} \\
        \midrule
        \multicolumn{3}{@{}l}{\textit{Total body variables}} \\
        $C_{\mathrm{insulin, blood}}$ & $66\,\mathrm{pM}$ & Blood insulin concentration \\
        $N_{\mathrm{cells, body}}$ & $2 \times 10^{12}$ & Total number of cells in the body \\
        $V_{\mathrm{blood}}$ & $5\,\mathrm{L}$ & Total blood volume \\
        $N_{\mathrm{insulin}}/N_{\mathrm{cells, body}}$ & $\sim 100$ & Insulin molecules per cell \\
        \multicolumn{3}{@{}l}{\textit{General system variables}} \\
        $R_{\mathrm{cell}}$ & $5\,\mu\mathrm{m}$ & Cell radius \\
        $R_{\mathrm{outer\,ring}}$ & $7.5\,\mu\mathrm{m}$ & Cell ring radius \\
        $R_{\mathrm{capillary}}$ & $5\,\mu\mathrm{m}$ & Capillary radius \\
        $L_{\mathrm{ring}}$ & $10\,\mu\mathrm{m}$ & Cylinder length \\
        $A_{\mathrm{ring}}$ & $471\,\mu\mathrm{m}^2$ & Ring surface area \\
        $N_{\mathrm{cells\,on\,ring}}$ & $\approx 4$ & Cells per ring \\
        $R_{\mathrm{receptor}}$ & $5\,\mathrm{nm}$ & Receptor radius \\
        $N_{\mathrm{receptors\,per\,cell}}$ & $10^{4}$ & Receptors per cell \\
        $\sigma_r$ & $84.9\,\mu\mathrm{m}^{-2}$ & Receptor surface density \\
        $D_{\mathrm{insulin}}$ & $125\,\mu\mathrm{m}^{2}/\mathrm{s}$ & Diffusion coefficient \\
        \multicolumn{3}{@{}l}{\textit{Insulin--receptor kinetics}} \\
        $K_{d, \mathrm{insulin}}$ & $2.7\,\mathrm{nM}$ & Dissociation constant \\
        $k_a$ & $3 \times 10^{6}\,\mathrm{M}^{-1}\mathrm{s}^{-1}$ & Association (binding) rate [CITE] \\
        $k_b$ & $4 \times 10^{-3}\,\mathrm{s}^{-1}$ & Dissociation (unbinding) rate [CITE] \\
        $k_c$ & $4 \times 10^{-3}\,\mathrm{s}^{-1}$ & Internalization rate [CITE] \\
        \bottomrule
    \end{tabular}
\end{table}

\section{Analytical Framework}\label{sec:analytical}
Incorporating physiological parameter values,
we can estimate the relative timescales of the system to infer if it is diffusion limited or not.
The goal of this is to see if the system is effectively at steady state regarding diffusion of ligands from the capillary wall to the outer cells,
relative to receptor timescales.
We model ligand sensing using a system of two coaxial cylinders.
The inner cylinder maintains a fixed ligand concentration, while the outer cylinder represents a cell membrane bearing discrete absorbing receptors embedded in an otherwise reflecting surface Fig.~\ref{fig:capillary_cells}.
We first analyze the idealized limit in which receptors act as perfect sinks, and subsequently incorporate receptor binding and unbinding kinetics as an additional resistance, rendering absorption finite.

\subsection{Separation of Time Scales}\label{sec:timescales}
The validity of the steady-state assumption, $\nabla^{2} C = 0$,
relies on a clear separation of the characteristic time scales governing diffusion, receptor kinetics, and systemic regulation.

The diffusive relaxation time
\begin{equation}
    \tau_{\mathrm{diff}} \approx \frac{(R-a)^{2}}{D} \approx 5 \times 10^{-2}\,\mathrm{s},
    \label{eq:tau_diff}
\end{equation}
corresponds to the rapid establishment of the insulin concentration profile between the capillary and the cell membrane.

The receptor kinetic timescale describes the equilibration of receptor occupancy with the local concentration $C(R)$,
\begin{equation}
    \tau_{\mathrm{kin}} = \frac{1}{k_a C(R) + k_b}.
    \label{eq:tau_kin}
\end{equation}
Using $k_a = 3 \times 10^{6}\,\mathrm{M}^{-1}\mathrm{s}^{-1}$, $k_b = 5 \times 10^{-3}\,\mathrm{s}^{-1}$, and physiological insulin levels, this yields $\tau_{\mathrm{kin}} \sim 10^{2}$--$10^{3}$\,s.
Finally, systemic insulin regulation occurs on a physiological timescale of several minutes, $\tau_{\mathrm{phys}} \sim 5$--$15$\,min.
Since $\tau_{\mathrm{diff}} \ll \tau_{\mathrm{kin}} \sim \tau_{\mathrm{phys}}$, diffusion is effectively instantaneous, while receptor kinetics evolve on the same scale as physiological insulin variations. The steady-state framework adopted here therefore applies to regimes in which insulin concentrations remain approximately constant for several minutes, such as basal conditions. Rapid transients may transiently violate this assumption and are not considered in this work.

\subsection{Geometric Resistance and Diffusive Capacity}\label{sec:geom_resistance}
To derive the steady-state flux, we adopt a diffusion-resistance analogy, in which the molecular flux $J$ is expressed as
\begin{equation}
    J = \frac{\Delta C}{\Rtot},
    \label{eq:flux_resistance}
\end{equation}
with $\Rtot$ the total effective diffusive resistance. Transport proceeds through two stages in series: diffusion from the source to the outer boundary, and capture by discrete receptors. The total resistance is therefore written as
\begin{equation}
    \Rtot = \Rgeom + \Rrec,
    \label{eq:rtot_geom_rec}
\end{equation}
where $\Rgeom$ depends only on the system geometry, and $\Rrec$ encodes receptor number and size.
For the cylinder-cylinder geometry, the geometric resistance is
\begin{equation}
    \Rgeom = \frac{\ln(R/a)}{2\pi H D},
    \label{eq:rgeom}
\end{equation}
while the resistance associated with $N$ absorbing receptors of radius $s$, acting in parallel, is
\begin{equation}
    \Rrec = \frac{1}{4 D N s}.
    \label{eq:rrec}
\end{equation}
Combining these resistances in series yields the total flux,
\begin{equation}
    J = \frac{\Jmax \, Ns}{\dfrac{\pi H}{2 \ln(R/a)} + Ns}
    \label{eq:flux_total_BP},
\end{equation}
with
\begin{equation}
    Jmax = \frac{2\pi H D C_{0}}{\ln(R/a)}
    \label{eq:flux_max}.
\end{equation}

This expression is formally identical to the classical Berg-Purcell result, with the sphere radius replaced by the effective geometric length scale $H/(2\ln(R/a))$.
All dependence on receptor number and size remains unchanged.
More generally, for a system containing $N$ disk-like perfectly absorbing receptors of radius $s$, the total steady-state flux can be written in the universal form
\begin{equation}
    J = \Jmax \, \frac{Ns}{Ns + \pi\,\mathrm{Cap(geom)}},
    \label{eq:flux_universal}
\end{equation}
where $\mathrm{Cap(geom)}$ is a purely geometric quantity characterizing the diffusive capacity of the system [CITE].
This formulation makes explicit that receptor number and size enter solely through the combination $Ns$, while all geometric information is encoded in the effective capacity.

\subsection{Finite Receptor Kinetics}\label{sec:finite_kinetics}
To move beyond the idealized assumption of perfectly absorbing receptors, we incorporate finite receptor binding kinetics. Within the diffusion--resistance framework, finite receptor kinetics introduce an additional resistance associated with ligand binding at the receptor surface. The total resistance is written as
\begin{equation}
    \Rtot = \Rdiff + \Rrxn,
    \label{eq:rtot_diff_rxn}
\end{equation}
where $\Rdiff$ is the diffusion resistance derived previously, and $\Rrxn$ is the kinetic resistance arising from receptor binding.

For $N$ identical receptors, each with intrinsic association rate $k_a$, the kinetic resistance is
\begin{equation}
    \Rrxn = \frac{1}{N k_a},
    \label{eq:rrxn}
\end{equation}
corresponding to $N$ identical resistors acting in parallel. Combining diffusion and reaction resistances in a series yields the total steady-state flux
\begin{equation}
    J = \frac{\Jmax \, Ns}{\alpha\pi + Ns + \dfrac{4\pi D \alpha s}{k_a}},
    \label{eq:flux_finite_kinetics}
\end{equation}
where $\alpha = H/(2\ln(R/a))$ is a geometric factor and $\Jmax = 2\pi H D C_{0}/\ln(R/a)$ is the diffusion-limited flux.
Equation~\eqref{eq:flux_finite_kinetics} reduces to the classical Berg-Purcell result in the limit $k_a \to \infty$, and correctly predicts vanishing flux as $k_a \to 0$.

\subsection{Homogenized boundary condition and effective surface reactivity}\label{sec:homogenized}
The substrate concentration $C(r,t)$ in the cylinder--cylinder geometry obeys the diffusion equation
\begin{equation}
    \frac{\partial C}{\partial t} = D\nabla^{2} C, \qquad a < r < R,
    \label{eq:diffusion_eq}
\end{equation}
with a fixed concentration imposed at the inner cylinder,
\begin{equation}
    C(a,t) = C_{0}.
    \label{eq:inner_bc}
\end{equation}

The outer cylinder at $r = R$ is sparsely covered by $N \gg 1$ small receptors of radius $s \ll R$, each of which absorbs ligands with finite intrinsic association rate $k_a$, while the remaining surface is reflecting.
Microscopically, this results in a mixed boundary condition with strong angular heterogeneity.
Our interest, however, lies in the \textit{total steady-state flux} into the receptor-covered boundary rather than in the local concentration field near individual receptors.
At steady state, flux conservation requires that the net radial flux through any cylindrical surface be identical.
This separation of scales allows us to replace the heterogeneous microscopic boundary by a homogenized, radially symmetric Robin boundary condition [CITE],
\begin{equation}
    -D \left.\frac{\partial C}{\partial r}\right|_{r=R} = \kastar \, \sigma_{0} \, C(R),
    \label{eq:robin_bc}
\end{equation}
where $\sigma_{0} = N/(2\pi R H)$ is the mean surface density of receptors and $\kastar$ is an effective surface reactivity.

\subsubsection{Derivation of the effective rate constant \texorpdfstring{$\kastar$}{ka*}}\label{sec:keff_derivation}
The effective reactivity $\kastar$ accounts for both intrinsic receptor kinetics and diffusion-limited access to the receptors.
It can be derived directly using a resistance-in-series argument.
Consider a single receptor of radius $s$.
In the limit of infinite intrinsic association rate ($k_a \to \infty$), the receptor acts as a perfect sink,
and the diffusion-limited capture rate is as follows
\begin{equation}
    k_{D} = 4 D s.
    \label{eq:kD}
\end{equation}

For finite $k_a$, diffusion toward the receptor and intrinsic binding at the receptor act as two resistances in series (Collins-Kimball theory).
The effective association rate for a single receptor is therefore
\begin{equation}
    \frac{1}{k_{\mathrm{eff}}} = \frac{1}{4Ds} + \frac{1}{k_a},
    \label{eq:keff_series}
\end{equation}
which yields
\begin{equation}
    k_{\mathrm{eff}} = 4Ds \left(1 + \frac{4Ds}{k_a}\right)^{-1}.
    \label{eq:keff_result}
\end{equation}

In the homogenized description, the total reaction flux is obtained by multiplying this single-receptor rate by the receptor surface density $\sigma_{0}$.
Matching to the boundary condition~\eqref{eq:robin_bc} therefore identifies
\begin{equation}
    \boxed{\kastar = 4Ds \left(1 + \frac{4Ds}{k_a}\right)^{-1}.}
    \label{eq:kastar}
\end{equation}

\subsubsection{Steady-state solution and flux}\label{sec:steady_state}
At steady state, the diffusion equation with boundary condition~\eqref{eq:robin_bc} admits the radially symmetric solution
\begin{equation}
    C(r) = C_{0}\left(1 - \mu \,\frac{\ln(r/a)}{\ln(R/a)}\right),
    \label{eq:C_profile}
\end{equation}
where the dimensionless uptake efficiency is
\begin{equation}
    \mu = \frac{\kastar \, \sigma_{0} \, R \, \ln(R/a)}{D + \kastar \, \sigma_{0} \, R \, \ln(R/a)}.
    \label{eq:mu}
\end{equation}

The total steady-state flux into the outer boundary is therefore
\begin{equation}
    J = \frac{2\pi H D C_{0}}{\ln(R/a)}\,\mu,
    \label{eq:J_mu}
\end{equation}
with
\begin{equation}
    \Jmax = \frac{2\pi H D C_{0}}{\ln(R/a)}
    \label{eq:Jmax_def}
\end{equation}
denoting the diffusion-limited flux corresponding to a perfectly absorbing outer cylinder.

\subsubsection{Physical interpretation}\label{sec:interpretation}
The effective rate constant $\kastar$ encapsulates the combined effects of receptor kinetics and diffusion limitation due to partial surface coverage.
In the perfect-sink limit $k_a \to \infty$, $\kastar$ reduces to $4Ds$, recovering the diffusion-limited Berg--Purcell result~\eqref{eq:flux_total_BP}.
For finite $k_a$, the uptake is reduced by the additional reaction resistance, yielding a smooth crossover between reaction-limited and diffusion-limited regimes~\eqref{eq:flux_finite_kinetics}.

\subsection{Receptor kinetics and Michaelis-Menten reduction at steady state}\label{sec:mm_reduction}
Following ligand binding, receptor--ligand complexes may either dissociate and release the ligand back into the environment or progress to a processed product. Receptor--ligand interactions are modeled using the standard substrate--enzyme reaction scheme
\begin{equation}
    S + E \mathrel{\mathop{\rightleftharpoons}^{k_a}_{k_b}} C \xrightarrow{k_c} P + E,
    \label{eq:mm_scheme}
\end{equation}
where $k_a$ is the association rate, $k_b$ the dissociation rate, and $k_c$ the processing (or internalization) rate [CITE].
Let $\sigma_f$, $\sigma_b$, and $\sigma_0$ denote the surface densities of free, bound, and total receptors, respectively.
Conservation of receptors implies
\begin{equation}
    \sigma_0 = \sigma_f + \sigma_b.
    \label{eq:receptor_conservation}
\end{equation}

\subsubsection{Steady-state Michaelis-Menten reduction}\label{sec:mm_ss}

In the long-time stationary regime, receptor occupancy and ligand concentration become time-independent [CITE]. Under these conditions, the surface reaction kinetics must satisfy
\begin{equation}
    \kastar \, C(R) \, \sigma_f = (k_b + k_c) \, \sigma_b,
    \label{eq:mm_ss_eq}
\end{equation}
where $\kastar$ is the effective association rate derived in the previous section, incorporating both intrinsic binding kinetics and diffusion-limited access to the receptors.

Solving Eq.~\eqref{eq:mm_ss_eq} together with receptor conservation yields the familiar Michaelis--Menten relations
\begin{equation}
    \sigma_f = \frac{\sigma_0 K_d}{C(R) + K_d}, \qquad \sigma_b = \frac{\sigma_0 C(R)}{C(R) + K_d},
    \label{eq:mm_relations}
\end{equation}
with effective dissociation constant
\begin{equation}
    K_d = \frac{k_b + k_c}{\kastar}.
    \label{eq:Kd_eff}
\end{equation}

Importantly, while the association rate $\kastar$ is diffusion-limited and geometry-dependent, the dissociation and processing rates $k_b$ and $k_c$ are local molecular processes and may therefore be taken directly from experimentally measured values.

\subsubsection{Flux balance at the receptor-bearing boundary}
\label{sec:flux_balance}

At steady state, the diffusion equation is supplemented by a flux balance condition at the receptor-bearing surface $r = R$, equating the diffusive flux to the net reaction flux,
\begin{equation}
    D \left.\frac{\partial C}{\partial r}\right|_{r=R} = \sigma_f \, \kastar \, C(R) - \sigma_b \, k_b.
    \label{eq:flux_balance}
\end{equation}

Substituting the Michaelis--Menten expressions~\eqref{eq:mm_relations} into Eq.~\eqref{eq:flux_balance} gives
\begin{equation}
    D \left.\frac{\partial C}{\partial r}\right|_{r=R} = \frac{\sigma_0 K_d}{C(R) + K_d}\,\kastar \, C(R) - \frac{\sigma_0 C(R)}{C(R) + K_d}\, k_b.
    \label{eq:flux_balance_mm}
\end{equation}

\subsubsection{Coupling to the steady-state concentration profile}
\label{sec:coupling}

The steady-state concentration profile in the cylindrical geometry is
\begin{equation}
    C(r) = C_{0}\left(1 - \mu \,\frac{\ln(r/a)}{\ln(R/a)}\right),
    \label{eq:C_profile_coupling}
\end{equation}
so that
\begin{equation}
    C(R) = C_{0}(1 - \mu), \qquad \left.\frac{\partial C}{\partial r}\right|_{r=R} = -\frac{C_{0}\mu}{R\ln(R/a)}.
    \label{eq:CR_dCdr}
\end{equation}

Substituting into the flux balance condition yields
\begin{equation}
    \frac{D C_{0}}{R \ln(R/a)}\,\mu = \frac{\sigma_0 K_d}{C_{0}(1-\mu) + K_d}\,\kastar C_{0}(1-\mu) - \frac{\sigma_0 C_{0}(1-\mu)}{C_{0}(1-\mu) + K_d}\,k_b.
    \label{eq:flux_balance_subst}
\end{equation}

Using $K_d = (k_b + k_c)/\kastar$, this expression simplifies to
\begin{equation}
    \frac{D}{R \ln(R/a)}\,\mu = \frac{\sigma_0 k_c (1-\mu)}{C_{0}(1-\mu) + K_d}.
    \label{eq:mu_implicit}
\end{equation}

\subsubsection{Solution and total flux}
\label{sec:solution_flux}

Equation~\eqref{eq:mu_implicit} yields a quadratic equation for $\mu$ with the physically admissible solution
\begin{equation}
    \mu = \frac{\dfrac{D(C_{0} + K_d)}{R \ln(R/a)} + \sigma_0 k_c - \sqrt{\left[\dfrac{D(C_{0} + K_d)}{R \ln(R/a)} + \sigma_0 k_c\right]^{2} - 4\dfrac{D C_{0}}{R \ln(R/a)} \sigma_0 k_c}}{2 \dfrac{D C_{0}}{R \ln(R/a)}}.
    \label{eq:mu_quadratic}
\end{equation}

We select the smaller root, ensuring $0 < \mu < 1$. In the absence of processing ($k_c \to 0$), this solution correctly yields $\mu \to 0$, corresponding to vanishing net uptake.

The total steady-state flux into the outer boundary is
\begin{equation}
    J = D \int_{r=R} \frac{\partial C}{\partial r}\,dS = \mu \Jmax,
    \label{eq:J_mu_final}
\end{equation}
where $\Jmax$ is the diffusion-limited (perfect sink) flux.

Since $\mu < 1$, finite receptor kinetics strictly reduce the uptake relative to the classical Berg--Purcell result,
\begin{equation}
    J < J_{\mathrm{BP}}.
    \label{eq:J_less_JBP}
\end{equation}

\subsubsection{Consistency with the Berg--Purcell limit}
\label{sec:consistency_BP}

As a consistency check, we recover the Berg--Purcell result in the perfect-kinetics limit. Taking $k_a \to \infty$ and $k_c \to \infty$ while keeping $k_b$ finite yields $K_d \to 0$ and $\mu \to 1$, recovering the diffusion-limited Berg--Purcell flux.

\section{Three-dimensional random-walk simulations}\label{sec:rw}
We performed three-dimensional random-walk (RW) simulations to model ligand diffusion in a coaxial cylinder--cylinder geometry.
The simulations were used both to validate the numerical scheme against known analytical results and to investigate systems with discrete absorbing patches on the outer cylinder surface, representing membrane receptors with finite kinetics.

\paragraph{Geometry.}
The system consists of an inner reflecting cylinder of radius $a$ and length $H$, coaxially enclosed by an outer cylinder of radius $R$ and the same length $H$.
Approximately $N \simeq 10^{4}$ particles are initialized uniformly within the annular volume $a \le r \le R$, $|z| \le H/2$.

\paragraph{Particle dynamics.}
Each particle undergoes Brownian motion with time step $\Delta t$, taking independent Gaussian-distributed displacements with standard deviation
\begin{equation}
    \sigma = \sqrt{2 D \Delta t},
    \label{eq:brownian_step}
\end{equation}
where $D$ is the diffusion coefficient.
The time step $\Delta t$ is chosen sufficiently small to resolve encounters with receptor patches and ensure numerical stability.

\paragraph{Boundary conditions.}
The inner cylinder ($r = a$) is reflecting and acts as a source that maintains a constant concentration $C_{0}$.
The outer cylinder ($r = R$) is reflecting except for discrete circular patches of radius $s$, which represent receptors.

In the \textit{steady-state} simulations, particles absorbed or processed at the outer boundary are immediately reintroduced at the inner cylinder to enforce flux conservation, $J_{R} = J_{a}$, thereby maintaining a stationary concentration profile (see Appendix~\ref{sec:appendix} for a proof of rapid convergence).

In the \textit{transient regime}, particles that reach an absorbing boundary are permanently removed from the system.

\paragraph{Receptor kinetics.}
When a diffusing particle intersects a receptor patch, binding occurs stochastically rather than deterministically.
Let $s$ denote the patch radius and $A_{\mathrm{patch}} = \pi s^{2}$ its area.
We define the surface reactivity $\kappa = k_a / A_{\mathrm{patch}}$, where $k_a$ is the intrinsic association rate.
Following Erban and Chapman [CITE], the probability of binding during a time step $\Delta t$ is
\begin{equation}
    p_{\mathrm{bind}} = \frac{\kappa \sqrt{\pi \Delta t / D}}{1 + \tfrac{1}{2} \kappa \sqrt{\pi \Delta t / D}}.
    \label{eq:p_bind}
\end{equation}

This expression accounts for short-time diffusive recollisions and ensures convergence to the correct macroscopic reaction rate as $\Delta t \to 0$.

Once bound, both the particle and the receptor are frozen.
The bound complex dissociates or is internalized with probabilities
\begin{equation}
    p_{\mathrm{unbind}} = 1 - e^{-k_b \Delta t}, \qquad p_{\mathrm{int}} = 1 - e^{-k_c \Delta t},
    \label{eq:p_unbind_int}
\end{equation}
corresponding to exponentially distributed waiting times with rates $k_b$ and $k_c$, respectively.
Upon dissociation, the particle is released a distance $O(s)$ from the patch on the outer cylinder, while internalized particles are reintroduced near the inner cylinder, modeling ligand processing and recycling Fig.~\ref{fig:rw_schematic}.

\begin{figure}[H]
    \centering
    \fbox{\parbox[c][7cm][c]{0.8\textwidth}{\centering \textit{Placeholder for Figure 2: Schematic of the random-walk simulation.}}}
    \caption{\textbf{Schematic of the random-walk simulation.} The annular domain is divided into four panels illustrating the key dynamical processes. \textbf{Top left:} diffusing particles undergoing specular reflection from the inner and outer boundaries. \textbf{Top right:} particle--receptor encounters at the outer boundary, resulting in either binding (absorption) or reflection with a finite binding probability. \textbf{Bottom left:} respawning of particles near the inner boundary following receptor-mediated internalization. \textbf{Bottom right:} respawning of particles near receptor sites after ligand dissociation (release).}
    \label{fig:rw_schematic}
\end{figure}

\subsection{Validation against analytical results}\label{sec:rw_validation}
All simulations in this section were performed with a fixed outer cylinder radius $R = 7.5$, diffusion coefficient $D = 1$, and cylinder length $H = 10$.
The integration time step was set to $\Delta t = 10^{-6}$.

\subsubsection{Steady-state flux to a fully absorbing cylinder .}\label{sec:rw_absorbing}

\paragraph{Radial concentration profile.}
We first verified the simulation by measuring the steady-state radial concentration profile.
The domain was partitioned into thin cylindrical shells, and the particle density was computed as a function of radius.
The measured profile closely followed the analytical solution of the diffusion equation, providing an additional consistency check Fig.~\ref{fig:rw_validation}a.

\paragraph{Particles absorption flux}
We further considered the simplest case in which the outer cylinder is perfectly absorbing.
Simulations were performed while varying the inner radius $a$, and the resulting steady-state flux was compared to the analytical Smoluchowski solution
\begin{equation}
    J_{\mathrm{cyl}} = \frac{2\pi H D C_{0}}{\ln(R/a)}.
    \label{eq:J_cyl}
\end{equation}

The numerical and analytical results showed good agreement Fig.~\ref{fig:rw_validation}b.
Small deviations were observed only for short simulation times or insufficient particle numbers and vanished upon increasing both, confirming the correctness of the implementation.

\subsubsection{Perfectly Absorbing Patches on the Outer Cylinder}\label{sec:rw_patches}
To probe more realistic receptor geometries, we next considered an outer cylinder that is reflecting except for $N$ small absorbing patches of radius $s$, uniformly distributed over its surface.
We focus on the regime where $s \ll R$.
In this limit, the normalized flux is expected to obey the Berg--Purcell equation~\eqref{eq:flux_total_BP}
\begin{equation}
    \frac{J}{\Jmax} = \frac{Ns}{Ns + \dfrac{\pi H}{2 \ln(R/a)}}, \qquad \Jmax = \frac{2\pi H D C_{0}}{\ln(R/a)}.
    \label{eq:J_patches_normalized}
\end{equation}

In these simulations, the receptor radius was fixed at $s = 0.1$, while the number of receptors $N$ was varied.
The normalized flux $J/\Jmax$ was computed and compared to the corresponding analytical prediction.
Across the explored parameter range, the numerical results show strong quantitative agreement with theory, demonstrating that the random-walk framework accurately captures diffusion-limited transport to discrete receptors in cylindrical geometries Fig.~\ref{fig:rw_validation}c.

\subsubsection{Transient state}\label{sec:rw_transient}
To validate the model and quantify the temporal evolution of the system, we monitored the total number of particles remaining in the domain as a function of time.
These results were compared with the analytical solution obtained numerically from Eq.~\eqref{eq:diffusion_eq}.
We examined the transient dynamics in two system configurations: (i) a fully absorbing outer cylinder, and (ii) a reflecting outer cylinder containing 5000 discrete perfect absorbing patches of radius $0.05$.
In both cases, the simulation results exhibit excellent quantitative agreement with the corresponding analytical predictions Fig.~\ref{fig:rw_validation}d.

\begin{figure}[H]
    \centering
    \fbox{\parbox[c][10cm][c]{0.85\textwidth}{\centering \textit{Placeholder for Figure 3: Validation of the cylinder-in-cylinder random-walk framework (4 panels: a, b, c, d).}}}
    \caption{\textbf{Validation of the cylinder-in-cylinder random-walk framework.} (A) Steady-state concentration profile for fully absorbing cylinder, Orange dots denote steady-state concentrations measured in shell segments from 3D RW simulations, while the solid blue line shows the analytical solution for a fully absorbing outer cylinder. (B) Steady-state flux for fully absorbing cylinder, expressed as an effective rate constant $k$, obtained from RW simulations (orange) compared with the analytical prediction (blue). (C) Steady-state flux of a system with patches on the outer cylinder compared for to the Berg--Purcell prediction for the cylindrical geometry~\eqref{eq:flux_total_BP}. (D) Transient survival probability for fully absorbing (Dirichlet) and partially absorbing (Robin) outer boundaries, showing slower decay for absorbing patches with excellent agreement between theory and simulation.}
    \label{fig:rw_validation}
\end{figure}

\subsubsection{Steady state with receptor kinetics}\label{sec:rw_kinetics}
After validating the random-walk simulation against analytical results for perfectly absorbing receptor patches, we progressively extended the model to incorporate finite receptor kinetics, as described above.
As in the analytical development, complexity was introduced incrementally, enabling systematic, stage-by-stage comparison between simulation and theory.

Validation was performed by comparing the normalized uptake flux $J/\Jmax$ (which is the effective absorption parameter $\mu$ calculated analytically in the simulations, with $\mu$ calculated analytically by Eqs.~\eqref{eq:mu_quadratic} and~\eqref{eq:J_mu_final}.

We first introduced a finite processing rate $k_c$, such that a particle captured by a receptor remains bound until release with probability $1 - \exp(-k_c \Delta t)$.
The resulting steady-state flux exhibited excellent agreement with the analytical predictions Fig.~\ref{fig:rw_kinetics}a.

Next, we incorporated a finite dissociation rate $k_b$, allowing bound particles to be released back into the medium at a distance of order the patch radius $s$ from the receptor.
This modification preserved quantitative agreement with the analytical solution, confirming that the release protocol accurately captures the effects of receptor unbinding Fig.~\ref{fig:rw_kinetics}b.

Finally, we introduced a finite association rate $k_a$, implemented via the stochastic binding rule described above.
Even in the fully kinetic regime with finite $k_a$, $k_b$, and $k_c$, the simulated uptake flux remained in excellent agreement with the analytical theory Fig.~\ref{fig:rw_kinetics}c.

Together, these results provide strong validation of both the analytical framework---specifically the homogenized boundary condition and the expression for the uptake efficiency $\mu$---and the stochastic simulation scheme.
Beyond validating the theory, the simulation framework offers additional flexibility for exploring regimes and mechanisms that are analytically intractable, making it a useful tool for future extensions of the model.

\begin{figure}[H]
    \centering
    \fbox{\parbox[c][12cm][c]{0.8\textwidth}{\centering \textit{Placeholder for Figure 4: Validation of the random walk simulation against the analytical theory for receptor kinetics uptake (3 panels: a, b, c).}}}
    \caption{\textbf{Validation of the random walk simulation against the analytical theory for receptor kinetics uptake.} Random-walk simulation results (symbols) are compared with analytical predictions from Eqs.~\eqref{eq:mu_quadratic} and~\eqref{eq:J_mu_final} (dashed lines). (a) Finite processing rate $k_c$ with perfectly absorbing receptors ($k_a \to \infty$, $k_b = 0$). (b) Addition of finite dissociation rate $k_b$, with particles released at a distance of order the receptor radius. (c) Fully kinetic receptors with finite association, dissociation, and processing rates. In all cases, excellent agreement is observed without adjustable parameters, validating both the homogenized boundary condition and the stochastic kinetic scheme.}
    \label{fig:rw_kinetics}
\end{figure}

\section{Physiological constraints on receptor-mediated ligand uptake}\label{sec:physiological}
Having validated the analytical model against RW simulation, we now use the theory to examine the behavior of the system under physiological conditions.
Specifically, we substitute representative physiological parameter values into the analytical expressions to assess how ligand uptake depends on receptor abundance and blood concentration $C_{0}$.

\paragraph{Flux efficiency.}
We focus on the normalized uptake efficiency $\mu = J/\Jmax$ as a function of receptor surface coverage for different ligand concentrations.
As shown in Fig.~\ref{fig:mu_coverage}, the uptake efficiency at physiological receptor numbers (on the order of $10^{5}$--$10^{6}$ receptors) is significantly below 1.
For example, at $N \sim 10^{5}$ receptors, $\mu \approx 5 \times 10^{-3}$, indicating that ligand uptake is far from diffusion-limited.

This behavior reflects the fact that insulin receptor kinetics are intrinsically slow, with characteristic rates on the order of $10^{-3}\,\mathrm{s}^{-1}$.
As a result, uptake is limited primarily by receptor kinetics rather than by diffusive transport to the cell surface.
In this regime, the steady-state concentration profile remains nearly flat, $C(r) \simeq C_{0}$, consistent with the limit $\mu \ll 1$ in Eq.~\eqref{eq:mu}.

Importantly, at low physiological ligand concentrations ($C_{0} \sim 1$--$100$\,pM), the uptake efficiency is higher and largely concentration-independent, reflecting unsaturated receptor kinetics.
At higher concentrations, receptor saturation becomes significant, reducing uptake efficiency as an increasing fraction of receptors remains occupied.

\begin{figure}[H]
    \centering
    \fbox{\parbox[c][7cm][c]{0.75\textwidth}{\centering \textit{Placeholder for Figure 5: $\mu = J/\Jmax$ vs.\ receptor surface coverage.}}}
    \caption{\textbf{Normalized ligand uptake flux.} The normalized flux $J/\Jmax$, computed from Eq.~\eqref{eq:mu_quadratic}, is shown as a function of receptor surface coverage on the outer cylinder for different ligand concentrations $C_{0}$. Physiological parameter values are used throughout.}
    \label{fig:mu_coverage}
\end{figure}

\paragraph{Receptor occupancy.}
Beyond the total ligand uptake flux, a key quantity of biological relevance is the receptor occupancy, since intracellular signaling is initiated by ligand-bound (and subsequently phosphorylated) receptors at the cell membrane.
In particular, downstream signaling pathways that regulate GLUT1 expression are driven by the fraction of occupied receptors rather than by the absolute ligand flux.

To examine this behavior, we compute the steady-state fraction of bound receptors as a function of the ligand concentration in the blood, $C_{0}$, for different total numbers of receptors (Fig.~\ref{fig:occupancy}).
One might expect that increasing the number of receptors would reduce the occupancy fraction, due to enhanced ligand consumption and a corresponding depletion of ligand near the cell surface.

Surprisingly, we find that the occupancy curves are nearly identical for different receptor abundances.
This behavior is a direct consequence of the fact that the uptake efficiency $\mu$ is much smaller than unity under physiological conditions.
As discussed above, when $\mu \ll 1$, ligand depletion at the cell surface is negligible and the local concentration at the receptors satisfies $C(R) \simeq C_{0}$.
Consequently, receptor occupancy is governed primarily by intrinsic receptor kinetics and depends only weakly on the total number of receptors.
For physiological insulin concentrations (10--500\,pM), the predicted receptor occupancy ranges from $\sim 0.3\%$ to 15\%.
Under basal conditions ($C_{0} \approx 50\,\mathrm{pM}$), occupancy is only $\sim 2\%$, indicating that receptors operate far from saturation.

\begin{figure}[H]
    \centering
    \fbox{\parbox[c][7cm][c]{0.75\textwidth}{\centering \textit{Placeholder for Figure 6: Bound receptor percentage vs $C_{0}$.}}}
    \caption{\textbf{Receptor occupancy under physiological conditions.} The steady-state fraction of bound receptors is computed from Eq.~\eqref{eq:mm_relations} as a function of the blood ligand concentration $C_{0}$ for different total receptor numbers. Because ligand depletion at the cell surface is negligible ($\mu \ll 1$), the occupancy curves collapse onto the Michaelis--Menten prediction and are largely independent of receptor abundance.}
    \label{fig:occupancy}
\end{figure}

\section{System constrains}\label{sec:constraints}
Now after we set the stage of our analytical equations of the system, we can start to get into conclusions about the system variables constrains.

\subsection{Constraint on Blood Insulin Concentration}\label{sec:blood_constraint}

\subsubsection{Lower bound on blood insulin concentration}\label{sec:lower_bound}
The minimum insulin concentration in the blood is constrained by the requirement that different cells reliably sense the hormone through receptor-mediated binding.
This reliability is limited by stochastic fluctuations in receptor occupancy across cells.

Receptor occupancy is described by the Michaelis--Menten relation,
\begin{equation}
    \theta = \frac{C(R)}{C(R) + K_d},
    \label{eq:theta_mm}
\end{equation}
where $\theta = \sigma_b / \sigma_0$ is the fraction of occupied receptors, $C(R)$ is the local insulin concentration at the cell surface, and $K_d$ is the effective dissociation constant defined in Eq.~\eqref{eq:Kd_eff}.

Assuming that the system is in steady state and that receptors act independently, each receptor undergoes a two-state binding--unbinding process with stationary occupancy probability $\theta$.
Consistent with the low receptor occupancy predicted under physiological basal conditions (typically $\theta \sim 2\%$), competition between receptors for ligand may be negligible, and receptor states may be treated as statistically independent.
Under these conditions, the instantaneous number of occupied receptors on a given cell follows a binomial distribution,
\begin{equation}
    N_b \sim \mathrm{Binomial}(\Nrec, \theta),
    \label{eq:binomial}
\end{equation}
where $\theta$ is given by the Michaelis--Menten relation.

The resulting coefficient of variation (CV) of receptor occupancy is
\begin{equation}
    \mathrm{CV} = \frac{\sqrt{\mathrm{Var}(N_b)}}{\langle N_b \rangle} = \sqrt{\frac{1 - \theta}{\Nrec \,\theta}}.
    \label{eq:CV}
\end{equation}

To ensure reliable sensing across cells, we impose the criterion $\mathrm{CV} < 0.1$, which yields the lower bound
\begin{equation}
    \theta > \frac{1}{1 + 0.01\,\Nrec}.
    \label{eq:theta_lower}
\end{equation}

Using the Michaelis--Menten equation, this condition translates into a constraint on the local insulin concentration,
\begin{equation}
    C(R) > \frac{100\,K_d}{\Nrec}.
    \label{eq:CR_lower}
\end{equation}

The local concentration $C(R)$ is related to the blood insulin concentration $C_{0}$ through the homogenized diffusion model derived above (Eq.~\eqref{eq:C_profile_coupling}), where the effective absorption parameter $\mu$ is determined by receptor kinetics and geometry (Eq.~\eqref{eq:mu_quadratic}).
Combining these results yields a lower bound on the blood insulin concentration,
\begin{equation}
    C_{0} > \frac{100\,K_d}{(1 - \mu)\,\Nrec}.
    \label{eq:C0_lower}
\end{equation}

Equation~\eqref{eq:C0_lower} provides a quantitative link between systemic insulin levels and the microscopic parameters governing receptor kinetics, spatial geometry, and sensing noise, and thus defines the minimal hormone concentration required for reliable cellular detection.

\subsubsection{Upper bound on blood insulin concentration}\label{sec:upper_bound}
An upper bound on the insulin concentration is imposed by the requirement that receptors do not become saturated.
At high occupancy, receptor-mediated sensing loses sensitivity, as the Michaelis--Menten response flattens and changes in ligand concentration no longer produce appreciable changes in receptor occupancy.

According to the Michaelis--Menten kinetics equation, sensitivity begins to decline once receptor occupancy exceeds approximately 50\%.
\begin{equation}
    \theta = \frac{C(R)}{C(R) + K_d} < \frac{1}{2},
    \label{eq:theta_upper}
\end{equation}
which implies
\begin{equation}
    C(R) < K_d.
    \label{eq:CR_upper}
\end{equation}

Relating the local surface concentration to the blood concentration via the diffusion model derived above, $C(R) = C_{0}(1 - \mu)$, we obtain the upper bound
\begin{equation}
    C_{0} < \frac{K_d}{1 - \mu}.
    \label{eq:C0_upper}
\end{equation}

Combining this constraint with the lower bound imposed by receptor occupancy noise (Eq.~\eqref{eq:C0_lower}), we find that the physiologically allowable insulin concentration lies within the window
\begin{equation}
    \frac{100\,K_d}{(1 - \mu)\,\Nrec} < C_{0} < \frac{K_d}{1 - \mu}.
    \label{eq:C0_window}
\end{equation}

Because $\mu$ depends on the ligand concentration $C_{0}$, we isolate $C_{0}$ from the inequality.
Doing so yields a $\mu$-independent upper bound on the admissible concentration,
\begin{equation}
    \frac{100\,K_d}{\Nrec} + \frac{100\,\sigma_0 k_c R \ln(R/a)}{D(\Nrec + 100)} < C_{0} < K_d + \frac{\sigma_0 k_c R \ln(R/a)}{2 D}.
    \label{eq:C0_mu_independent}
\end{equation}

Equation~\eqref{eq:C0_mu_independent} defines a finite dynamic range for systemic insulin concentration, determined jointly by receptor quantity, kinetics, transport geometry, and the requirement of reliable cellular sensing.

\subsection{Constraints on Receptor Number}\label{sec:receptor_constraints}
When formulating constraints on receptor abundance, we seek bounds that are \textit{independent} of the ambient insulin concentration $C_{0}$.
The goal is to choose a receptor number that remains relevant across conditions (ensuring reliable capture), while avoiding excessive membrane crowding and steric interference that would violate the independent-receptor assumption.

\subsubsection{Minimum constraint}\label{sec:min_constraint}
A natural lower bound is obtained by requiring that uptake operates at or above a prescribed fraction of the maximal possible flux,
\begin{equation}
    \frac{J}{\Jmax} \ge \left(\frac{J}{\Jmax}\right)_{\min}, \qquad 0 < \left(\frac{J}{\Jmax}\right)_{\min} < 1.
    \label{eq:flux_min}
\end{equation}

In principle, one could enforce the constraint~\eqref{eq:flux_min} using the full finite-kinetics expression for the flux (Eq.~\eqref{eq:J_mu_final}).
However, this would yield a receptor-number constraint that depends explicitly on the ligand concentration $C_{0}$ and on receptor kinetic parameters.
To obtain a conservative and \textit{concentration-independent} bound, we instead evaluate~\eqref{eq:flux_min} in the diffusion-limited (Berg--Purcell) regime, Eq.~\eqref{eq:flux_total_BP}.

This choice isolates the purely diffusive accessibility of receptors and provides a minimal requirement on receptor number that is independent of binding kinetics.
Physically, this bound ensures that ligand molecules can reach the receptor-bearing surface at a sufficient rate, irrespective of downstream reaction processes.
Incorporating finite receptor kinetics can only reduce the uptake efficiency and therefore increase the required number of receptors.

Substituting the Berg--Purcell form and solving for $N$ yields the conservative bound
\begin{equation}
    N \;\ge\; \frac{\left(\dfrac{J}{\Jmax}\right)_{\min}}{1 - \left(\dfrac{J}{\Jmax}\right)_{\min}}\,\frac{\pi H}{2 s \ln(R/a)}, \qquad 0 < \left(\frac{J}{\Jmax}\right)_{\min} < 1.
    \label{eq:N_min}
\end{equation}

\subsubsection{Upper bound on receptor density}\label{sec:receptor_upper}
An upper bound on receptor abundance is set by steric crowding on the cell membrane.
Let $\varepsilon$ denote the typical spacing between neighboring receptors.
When this spacing becomes comparable to twice the receptor radius, $\varepsilon \lesssim 2 R_{\mathrm{rec}}$, receptors physically contact each other, lose lateral mobility, and the independent-receptor approximation breaks down.

Approximating receptors as uniformly distributed over the cylinder surface area $\Acyl$, the mean spacing scales as
\begin{equation}
    \varepsilon \sim \sqrt{\frac{\Acyl}{\Nrec}}.
    \label{eq:epsilon}
\end{equation}

Requiring $\varepsilon > 2 R_{\mathrm{rec}}$ yields the steric constraint
\begin{equation}
    \Nrec < \frac{\Acyl}{4 R_{\mathrm{rec}}^{2}}.
    \label{eq:N_steric}
\end{equation}

Importantly, receptors do not occupy the membrane in isolation.
The human genome encodes on the order of $10^{3}$--$10^{4}$ surface proteins involved in cell--cell and cell--environment interactions, with a typical cell expressing a few hundred distinct membrane protein types [CITE].
As a rough but conservative estimate, we take
\begin{equation}
    \Nsurf \simeq 100
    \label{eq:Nsurf}
\end{equation}
distinct surface protein species sharing the available membrane area.
Assuming comparable surface coverage among these species, the maximal allowable number of receptors of a given type is reduced to
\begin{equation}
    \Nrec < \frac{\Acyl}{4 \Nsurf R_{\mathrm{rec}}^{2}} = \frac{\Acyl}{1200\,R_{\mathrm{rec}}^{2}}.
    \label{eq:N_upper}
\end{equation}

This bound is necessarily approximate, but it provides a physically motivated upper limit on receptor abundance that is independent of ligand concentration.

Collecting the steric upper bound and the functional lower bound, the admissible receptor number is constrained to the interval
\begin{equation}
    \frac{\left(\dfrac{J}{\Jmax}\right)_{\min}}{1 - \left(\dfrac{J}{\Jmax}\right)_{\min}}\,\frac{\pi H}{2 s \ln(R/a)} \;\le\; N \;<\; \frac{\Acyl}{1200\,R_{\mathrm{rec}}^{2}}.
    \label{eq:N_window}
\end{equation}

\subsubsection{Numerical values}\label{sec:numerical_values}
Substituting representative physiological parameters into Eq.~\eqref{eq:N_window} and requiring a minimum relative flux $(J/\Jmax)_{\min} = 0.7$, we obtain the following numerical design window for the number of receptors:
\begin{equation}
    1.8 \times 10^{4} \;\lesssim\; N \;\lesssim\; 4.7 \times 10^{4}.
    \label{eq:N_numerical}
\end{equation}

This range suggests a characteristic receptor abundance of order $N \sim 3 \times 10^{4}$.

Using this representative value, together with the physiological parameters listed in Table~\ref{tab:parameters}, in the concentration bounds derived above, we obtain the corresponding admissible window for the insulin concentration,
\begin{equation}
    14.75\,\mathrm{pM} \;\lesssim\; C_{0} \;\lesssim\; 2.7\,\mathrm{nM}.
    \label{eq:C0_numerical}
\end{equation}

Notably, the lower bound is of the same order of magnitude as the lower end of physiological insulin levels in healthy individuals, while the upper bound is of the same order as concentrations reported under hyperinsulinemic conditions [CITE].

\section{Conclusions}\label{sec:conclusions}
In this work, we derived an analytical expression for ligand uptake that simultaneously incorporates receptor kinetics ($k_a$, $k_b$, $k_c$), system geometry, and receptor abundance (Eq.~\eqref{eq:mu_quadratic}), and validated the theory using three-dimensional random-walk simulations.

A central finding of our analysis is that, under physiological conditions, insulin uptake is \textit{kinetically limited} rather than diffusion-limited.
In the classical Berg--Purcell regime, sufficiently abundant receptors render uptake insensitive to receptor number, making variability in receptor expression physiologically irrelevant.
However, insulin receptors are expressed over a wide dynamic range across tissues and disease states, implying that uptake must operate outside this diffusion-limited regime.

Slow dissociation and internalization rates ($k_b$, $k_c$) ensure that once insulin binds, receptors remain occupied for sufficiently long times to initiate downstream signaling.
At the same time, the association rate ($k_a$) must remain moderate: if binding were too fast, receptors would saturate even at basal insulin levels, eliminating sensitivity to changes in hormone concentration.
Together, these constraints enforce a kinetically controlled uptake regime in which flux scales with both ligand concentration and receptor abundance.

This kinetic balance preserves insulin as a graded, information-rich endocrine signal rather than a rapidly absorbed metabolite.
Sensitivity is maintained across physiological concentrations, while receptor saturation and diffusion-limited behavior are avoided, allowing cells to tune their response through both receptor number and ligand availability..

In the physiological regime considered here, receptor occupancy is low, with only $\sim 2\%$ of receptors bound at basal insulin concentrations.
In this limit, the Michaelis--Menten relation reduces to the linear form
\begin{equation}
    \theta \simeq \frac{C\,\Nrec}{K_d},
    \label{eq:theta_linear}
\end{equation}
so that the number of occupied receptors scales linearly with both insulin concentration and receptor abundance.
At first glance, this appears problematic: receptor expression varies substantially across cells, often by 20\% or more, suggesting comparable variability in the number of activated receptors.
However, this variability may be biologically inconsequential.
For example, a 20\% change corresponds to differences on the order of 200 versus 240 occupied receptors, which may fall well below the resolution threshold of downstream signaling pathways.

By contrast, postprandial insulin levels rise by a factor of 5--10, producing changes in receptor activation that are an order of magnitude larger and therefore readily distinguishable by the cell.
How such differences in receptor activation are integrated and transduced into GLUT transporter production depends on intracellular signaling networks and transcriptional regulation, which lie beyond the scope of the present work.
Nonetheless, our analysis clarifies the physical and kinetic constraints under which insulin signaling operates, and highlights the role of receptor kinetics in shaping both sensitivity and robustness of hormonal sensing.

Notably, insulin receptor abundance differs by more than an order of magnitude across tissues, ranging from $\sim 10^{4}$ receptors per cell in skeletal muscle to $\sim 10^{6}$ in adipocytes, with hepatocytes occupying an intermediate range [CITE, CITE, CITE].
Such large differences strongly suggest that insulin signaling is not designed to resolve small ($\sim 10$--$20\%$) variations in receptor number.
Instead, meaningful changes in cellular response must correspond to order-of-magnitude differences in receptor activation or ligand concentration.

Moreover, if insulin capture were saturated at low receptor numbers or at basal hormone levels, all cell types---regardless of their receptor abundance---would experience similar receptor occupancy and activation.
This would erase tissue-specific sensitivity and contradict the well-established heterogeneity of insulin responses across tissues.
The observed order-of-magnitude variation in receptor number therefore provides strong evidence that insulin uptake operates in a kinetically restricted regime, in which receptor occupancy remains low and scales with both ligand concentration and receptor abundance.

Regarding system limitations, we examined constraints on the number of receptors by considering both physical crowding at the cell surface as an upper bound and diffusion-limited efficiency as a lower bound, as described by Eq.~\eqref{eq:flux_total_BP}.
This analysis yields a permissible range of receptor numbers that is consistent with biologically observed values Eq.~\eqref{eq:N_window}.

We further investigated substrate concentration constraints using the derived receptor-occupancy equilibrium equation, Eq.~\eqref{eq:mu_quadratic}, in the Michaelis--Menten regime.
By defining a lower bound based on a 10\,\% coefficient of variation between cells, assuming identical amount of receptors between cells and independent binding events of receptros.\ and an upper bound based on receptor occupancy sensitivity, we obtained a concentration range whose lower limit closely matches physiological levels.
However, the upper bound predicted by the model remains significantly higher, reaching the nanomolar range, whereas physiological concentrations are typically on the order of hundreds of picomolar Eq.~\eqref{eq:C0_mu_independent}.
This discrepancy indicates that the present model does not yet capture the full complexity of the biological system.


\appendix
\section*{Appendix}
\addcontentsline{toc}{section}{Appendix}
\section{Rapid relaxation induced by particle respawning at the system boundary}\label{sec:appendix}
In the stochastic simulations, particles absorbed at the outer boundary are respawned at the inner boundary in order to enforce flux conservation and maintain a stationary concentration profile.
Here, we provide a simple analytical argument demonstrating that this respawning mechanism not only produces the correct steady state, but also accelerates relaxation toward it.

To this end, we consider a one-dimensional diffusion problem on the interval $0 \le x \le L$,
\begin{equation}
    \frac{\partial c}{\partial t} = D \frac{\partial^{2} c}{\partial x^{2}}, \qquad c(0,t) = 0,
    \label{eq:1D_diff}
\end{equation}
and decompose the concentration into a steady-state profile and a transient deviation,
\begin{equation}
    c(x,t) = c_{\mathrm{st}}(x) + u(x,t).
    \label{eq:decomposition}
\end{equation}

The deviation $u(x,t)$ satisfies
\begin{equation}
    \frac{\partial u}{\partial t} = D \frac{\partial^{2} u}{\partial x^{2}}, \qquad u(0,t) = 0,
    \label{eq:u_eq}
\end{equation}
together with a boundary condition at $x = L$ that depends on how absorbed particles are handled.

\subsection{Case 1: Dirichlet boundary condition}\label{sec:case1}
If the concentration is fixed at the outer boundary, $c(L,t) = c_{0}$, then the deviation satisfies
\begin{equation}
    u(L,t) = 0.
    \label{eq:bc_dirichlet}
\end{equation}

Using separation of variables, $u(x,t) = X(x) T(t)$, the spatial eigenvalue problem becomes
\begin{equation}
    X'' + k^{2} X = 0, \qquad X(0) = X(L) = 0,
    \label{eq:eigen_dirichlet}
\end{equation}
with eigenfunctions $X_{n}(x) = \sin(k_{n} x)$ and eigenvalues
\begin{equation}
    k_{n} = \frac{n\pi}{L}, \qquad n = 1, 2, \ldots
    \label{eq:kn_dirichlet}
\end{equation}

The slowest decaying mode corresponds to $n = 1$, yielding a relaxation time
\begin{equation}
    \tau_{1} = \frac{1}{D(\pi/L)^{2}} = \frac{L^{2}}{\pi^{2} D}.
    \label{eq:tau1}
\end{equation}

\subsection{Case 2: Slope-matching boundary condition}\label{sec:case2}
In the respawning scheme used in the simulations, flux conservation is enforced by requiring that the diffusive flux entering the system at $x = 0$ matches the flux leaving at $x = L$.
This corresponds to a slope-matching condition,
\begin{equation}
    \left.\frac{\partial c}{\partial x}\right|_{x=0} = \left.\frac{\partial c}{\partial x}\right|_{x=L},
    \label{eq:slope_match}
\end{equation}
which implies for the deviation
\begin{equation}
    \left.\frac{\partial u}{\partial x}\right|_{x=0} = \left.\frac{\partial u}{\partial x}\right|_{x=L}.
    \label{eq:slope_match_u}
\end{equation}

The spatial eigenvalue problem now becomes
\begin{equation}
    X'' + k^{2} X = 0, \qquad X(0) = 0, \qquad X'(0) = X'(L),
    \label{eq:eigen_slope}
\end{equation}
with solutions $X_{n}(x) = \sin(k_{n} x)$ and eigenvalues determined by
\begin{equation}
    \cos(k_{n} L) = 1 \implies k_{n} = \frac{2\pi n}{L}, \qquad n = 1, 2, \ldots
    \label{eq:kn_slope}
\end{equation}

The slowest decaying mode corresponds to $n = 1$, yielding a relaxation time
\begin{equation}
    \tau_{2} = \frac{1}{D(2\pi/L)^{2}} = \frac{L^{2}}{4\pi^{2} D}.
    \label{eq:tau2}
\end{equation}

\subsection{Comparison and implication}\label{sec:comparison}
Comparing the two relaxation times,
\begin{equation}
    \tau_{2} = \frac{1}{4}\tau_{1},
    \label{eq:tau_ratio}
\end{equation}
we find that the slope-matching (respawning) boundary condition leads to relaxation that is four times faster than the standard Dirichlet boundary condition.
This acceleration arises because the lowest admissible spatial mode has twice the wavenumber, resulting in steeper gradients and faster diffusive decay.

This analysis demonstrates that respawning absorbed particles at the system boundary not only preserves the correct steady-state concentration profile, but also significantly enhances convergence toward that steady state.
As a result, the respawning scheme provides an efficient and well-controlled method for enforcing stationarity in stochastic diffusion simulations.


\begin{thebibliography}{99}

\bibitem{wananant2000} Sumanas Wananant and Michael J. Quon. Insulin receptor binding kinetics: Modeling and simulation studies. \textit{Journal of Theoretical Biology}, 205(2):331--346, 2000.

\bibitem{marshall1985} Stephen Marshall. Kinetics of insulin receptor internalization and recycling in adipocytes. \textit{Journal of Biological Chemistry}, 260(7):4136--4144, 1985.

\bibitem{berg1977} Howard C. Berg and Edward M. Purcell. Physics of chemoreception. \textit{Biophysical Journal}, 20(2):193--219, 1977.

\bibitem{ward1993} Michael J. Ward and Joseph B. Keller. Strong localized perturbations of elliptic problems. \textit{SIAM Journal on Applied Mathematics}, 53(3):770--798, 1993.

\bibitem{handy2021} Gregory Handy and Sean D. Lawley. Revising Berg--Purcell for finite receptor kinetics. \textit{Biophysical Journal}, 120(11):2237--2248, 2021.

\bibitem{gunawardena2007} Jeremy Gunawardena. Distributivity and processivity in multisite phosphorylation can be distinguished through steady-state invariants. \textit{Biophysical Journal}, 93(11):3828--3834, 2007.

\bibitem{keener2010} James P. Keener and James Sneyd. \textit{Mathematical Physiology: I. Cellular Physiology}. Springer, New York, 2010.

\bibitem{erban2007} Radek Erban and S. Jonathan Chapman. Reactive boundary conditions for stochastic simulations of reaction-diffusion processes. \textit{Physical Biology}, 4(1):16, 2007.

\bibitem{uhlen2024} Mathias Uhl\'en, Markus J. Karlsson, Wenzhong Zhong, Abdellah Tebani, et al. A genome-wide map of human surface proteins, 2024.

\bibitem{hiriart2014} M. Hiriart et al. Hyperinsulinemia is associated with increased soluble insulin receptor in plasma of obese subjects. \textit{Frontiers in Endocrinology}, 5:95, 2014. This study uses 1\,nmol\,L$^{-1}$ insulin exposure in hepatocyte experiments to model high physiological concentrations that occur during hyperinsulinemia.

\bibitem{pringault1983} E. Pringault and C. Plas. Differences in degradation processes for insulin and its receptor in cultured foetal hepatocytes. \textit{Biochemical Journal}, 214(1):191--198, 1983.

\bibitem{im1986} M. J. Im et al. Insulin binding sites per cell in isolated myocytes. \textit{Biochemical and Biophysical Research Communications}, 135(2):879--885, 1986.

\bibitem{pagano1977} G. Pagano et al. Insulin receptors in adipocytes of non-diabetic and diabetic subjects. \textit{Acta Diabetologica Latina}, 14(2):98--106, 1977.

\end{thebibliography}


\end{document}
